%-------------------------
%---- Paquetes requeridos ----
\documentclass[a4paper,11pt]{article}

\usepackage{import}
\subimport{}{packages.tex}
\subimport{}{macros.tex}

\usepackage[utf8]{inputenc}
\usepackage{multicol}
\usepackage{ragged2e}

%-------------------------------------------
%%%%%% INICIO DEL CV %%%%%%
%%%%%% DEFINICIÓN DE DATOS %%%%%%

\newcommand{\name}{Jeanpier Alexander Robles Fabian}
\newcommand{\phone}{926 776 359}
\newcommand{\emaila}{jeanpierrf31@gmail.com}
\newcommand{\linkedin}{https://www.linkedin.com/in/jeanpier-robles/}
\newcommand{\github}{https://github.com/jeannrf}
\newcommand{\location}{Lima, Perú}

\begin{document}
\fontfamily{cmr}\selectfont

%----------ENCABEZADO-----------------
\begin{center}
    {\Huge\bfseries \name\par}



    \vspace{2mm}

    {\small
        \location
        \enspace\textbar\enspace
        \href{mailto:\emaila}{\emaila}
        \enspace\textbar\enspace
        +51 \phone
        \par
    }

    \vspace{1mm}

    {\small
        \href{\linkedin}{LinkedIn}
        \enspace\textbar\enspace
        \href{\github}{GitHub}
        \par
    }
\end{center}

\vspace{2mm}


%-------------PROFILE-----------------
\begin{addmargin}[\leftM]{\rightM}
\justifying
Estudiante de Ingeniería de Software de la UNI, perteneciente al décimo superior, con experiencia en proyectos de desarrollo web, ciencia de datos e inteligencia artificial.
He participado en la creación de soluciones aplicadas a desarrollo de sistemas, automatización y gestión de información.
Cuento con experiencia liderando equipos de desarrollo y coordinando proyectos tecnológicos multidisciplinarios.

\end{addmargin}

%-----------EDUCATION-----------------
\section{\textbf{FORMACIÓN ACADÉMICA}}
\begin{tabularx}{\textwidth}{p{14cm} R}
  \textbf{Estudiante de Ingeniería de Software - Quinto Ciclo} & \textit{\footnotesize Agosto 2024 -- Actualidad} \\
  Universidad Nacional de Ingeniería (UNI) & \\
  \small Décimo Superior - Top 7 de la Especialidad (Ciclo 26-1) &
\end{tabularx}


%-----------LEADERSHIP-----------------
\section{\textbf{LIDERAZGO Y ACTIVIDADES EXTRACURRICULARES}}

%--- Liderazgo y Actividad 1 ---
\noindent
\begin{tabularx}{\textwidth}{@{}X r@{}}
\textbf{Líder organizador - Software Week 2026} & \textit{\small Junio 2026 -- Septiembre 2026}
\end{tabularx}

\begin{itemize}[leftmargin=1.5em, nosep]
\item Coordinación con 4 áreas y planificación del congreso, cronograma de marketing, creación de bases de las actividades, negociación con patrocinadores, desarrollo de plataforma web con gestión de base de datos.
\end{itemize}

\vspace{3mm}

%--- Liderazgo y Actividad 2 ---
\noindent
\begin{tabularx}{\textwidth}{@{}X r@{}}
\textbf{Director del área de IT - TinkuyLabs} & \textit{\small Marzo 2026 -- Presente}
\end{tabularx}

\begin{itemize}[leftmargin=1.5em, nosep]
\item Planificación de proyectos tecnológicos como página web, sistema de inscripción, entregables via web.
\end{itemize}

\vspace{3mm}

%--- Liderazgo y Actividad 3 ---
\noindent
\begin{tabularx}{\textwidth}{@{}X r@{}}
\textbf{Director de Relaciones Públicas - Unicode} & \textit{\small Agosto 2025 -- Agosto 2026}
\end{tabularx}

\begin{itemize}[leftmargin=1.5em, nosep]
\item Realización de 4 eventos tecnológicos de alto impacto llegando a tener picos de +180 asistentes dentro de la UNI, gestión de alianzas, coordinación y contacto con ponentes del sector tech. 
\end{itemize}

\vspace{2mm}

%-----------------Habilidades & Certificaciones-----------


%-----------PROJECTS-----------------
\section{\textbf{PROYECTOS Y EXPERIENCIA RELEVANTE}}

%--- Project 1: Cupo Cero ---
\noindent
\begin{tabularx}{\textwidth}{@{}X r@{}}
\textbf{Cupo Cero — Multiplataforma de reporte de extorsiones} & \textit{\small Agosto 2026}
\end{tabularx}
\small
Sistema móvil y web para la geolocalización de zonas de riesgo por extorsión nacional, diseñado para apoyar a las municipalidades y agilizar el proceso de denuncia de las víctimas.

\begin{itemize}[leftmargin=1.5em, nosep]
\item Integré Gemma con LiteRT-LM para extraer datos de chats en el dispositivo, preservando la privacidad.
\item Desarrollé un dashboard con mapas de calor, trazabilidad de denuncias e integración con la app móvil.

\end{itemize}

\textbf {Tecnologías:} React, Gemma LLM, Google AI Edge (LiteRT-LM), Firebase, Cloud Run, Firestore

\vspace{5mm}

%--- Project 2: Unicode Landing Page ---
\noindent
\begin{tabularx}{\textwidth}{@{}X r@{}}
\textbf{Unicode Page -- Plataforma Institucional} & \textit{\small Marzo 2026 -- Abril 2026}
\end{tabularx}

\small
Plataforma web desarrollada para la visibilidad de la organización, impactando a una comunidad y gestionando la difusión de proyectos y eventos tecnológicos. En este proyecto tomé el rol de líder de un equipo de 4 personas.

\begin{itemize}[leftmargin=1.5em, nosep]
\item Lideré a 4 estudiantes de Software usando Git/GitHub, revisión de pull requests y desarrollo con buenas prácticas.
\item Implementé el backend de convocatorias con PostgreSQL, postulaciones, autenticación y control de acceso.
\end{itemize}

\textbf{Tecnologías:} Astro, TypeScript, Tailwind CSS, Supabase (PostgreSQL), Vercel, Git, GitHub

\vspace{5mm}
%--- Project 3: AutoML Advisor ---
\noindent
\begin{tabularx}{\textwidth}{@{}X r@{}}
\textbf{AutoML Advisor – Consultor Inteligente de Ciencia de Datos} & \textit{\small Diciembre 2025}
\end{tabularx}

\small
Aplicación web inteligente que automatiza el análisis de datasets y traduce resultados de machine learning en recomendaciones comprensibles para el usuario.
\begin{itemize}[leftmargin=1.5em, nosep]
    \item Implementé preprocesamiento de datos, división 80/20 de entrenamiento y prueba comparando modelos de Regresión Logística, Regresión Lineal y Random Forest para problemas de clasificación y regresión.
    \item Integré Groq API para interpretar métricas, detectar problemas de calidad y generar recomendaciones.
\end{itemize}

\textbf{Tecnologías:} Python, FastAPI, Pandas, NumPy, Scikit-learn, Groq API, JavaScript, HTML, CSS.

\vspace{5mm}

%-----------------Habilidades & Certificaciones-----------
\section{\textbf{HABILIDADES Y CERTIFICACIONES}}

\begin{itemize}[leftmargin=1.5em, nosep]
\item \textbf{Desarrollo Frontend:} Astro, Tailwind CSS, TypeScript, React, JavaScript, HTML, CSS.
    \item \textbf{Backend y Datos:} Python (Scikit-learn, Pandas), FastAPI, Java, PostgreSQL, Supabase, Firebase.
    \item \textbf{IA y Productividad:} Herramientas de IA para desarrollo de software, Figma, Ofimática.
    \item \textbf{Certificaciones:} Introduction to DevOps (IBM), Introducción al Desarrollo de Software (PUCP), Generative AI: Introduction and Applications (IBM), Introduction to Data Science (IBM) Liderazgo y Comunicación (U. Palermo).
    \item \textbf{Idiomas:} Español (Nativo), Inglés (B1)
\end{itemize}

\end{document}
